% eksperymentalne tłumaczenie na włoski

%

\subsection{Aksjomaty incydencji} % lang-pl
\index{aksjomat!incydencji}% % lang-pl
\subsection{Assiomi di incidenza} % lang-it
\index{assioma!di incidenza}% % lang-it

\begin{axiom}
[incydencji, I1] % lang-pl
[di incidenza, I1] % lang-it
    Dla każdej pary punktów $A$ oraz $B$ istnieje dokładnie jedna prosta $l$, która przechodzi przez te punkty. % lang-pl
    Per ogni coppia di punti $A$ e $B$ esiste un'unica retta $l$ che passa per tali punti. % lang-it
\end{axiom}

Zdania ,,prosta $l$ przechodzi przez punkt $A$'' oraz ,,punkt $A$ leży na prostej $l$'' mają dokładnie to samo znaczenie i będziemy je stosować zamiennie. % lang-pl
Le espressioni ``la retta $l$ passa per il punto $A$'' e ``il punto $A$ appartiene alla retta $l$'' hanno esattamente lo stesso significato e le useremo come sinonimi. % lang-it

\begin{axiom}
[incydencji, I2] % lang-pl
[di incidenza, I2] % lang-it
    Na każdej prostej leżą co najmniej dwa punkty. % lang-pl
    Su ogni retta giacciono almeno due punti. % lang-it
\end{axiom}

\begin{axiom}
[incydencji, I3] % lang-pl
[di incidenza, I3] % lang-it
    Istnieją co najmniej trzy punkty, które nie są współliniowe. % lang-pl
    Esistono almeno tre punti che non sono allineati. % lang-it
\end{axiom}

,,Być współliniowym'' to synonim leżenia na jednej prostej. % lang-pl
\index{współliniowy}% % lang-pl
Analogicznie mówimy, że wiele prostych, które przechodzą przez jeden punkt, jest współpękowych. % lang-pl
``Essere allineati'' significa giacere sulla stessa retta. % lang-it
\index{allineato}% % lang-it
Analogamente, diciamo che più rette che passano per uno stesso punto sono concorrenti. % lang-it

\begin{proposition}
    Dwie różne proste mogą mieć co najwyżej jeden punkt wspólny. % lang-pl
    Due rette distinte possono avere al più un punto in comune. % lang-it
\end{proposition}

\begin{definition}
    Dwie różne proste, które nie są współpękowe (nie mają punktów wspólnych), nazywamy równoległymi. % lang-pl
    Każda prosta jest też równoległa do siebie. % lang-pl
    Due rette distinte che non sono concorrenti (non hanno punti in comune) si dicono parallele. % lang-it
    Ogni retta è inoltre parallela a sé stessa. % lang-it
\end{definition}
\index{prosta!równoległa}% % lang-pl
\index{retta!parallela}% % lang-it

Symbolicznie: $k \parallel l$. % lang-pl
Trzeba uważać, bo dla niektórych dwie proste równoległe są współpękowe, bo przechodzą przez punkt w nieskończoności, jak zwykło się mawiać w geometrii rzutowej. % lang-pl
Simbolicamente: $k \parallel l$. % lang-it
Bisogna fare attenzione, perché per alcuni due rette parallele sono concorrenti, poiché passano per un punto all'infinito, come si usa dire nella geometria proiettiva. % lang-it

Na dwóch prostych, które nie są równoległe ($k \nparallel l$), leży dokładnie jeden punkt. % lang-pl
Mówimy, że proste się przecinają, a wspomniany punkt jest ich punktem przecięcia. % lang-pl
Oznaczamy go $k \cdot l$. % lang-pl
Su due rette che non sono parallele ($k \nparallel l$) giace esattamente un punto. % lang-it
Diciamo che le rette si intersecano e che il punto menzionato è il loro punto di intersezione. % lang-it
Lo indichiamo con $k \cdot l$. % lang-it

Hartshorne \cite[s. 38]{hartshorne_2000} przywoła: % lang-pl
Hartshorne \cite[p. 38]{hartshorne_2000} presenta il seguente assioma: % lang-it

\begin{axiom}
[Playfaira, P] % lang-pl
[di Playfair, P] % lang-it
    Dla każdej prostej $l$ oraz punktu $A$, istnieje co najwyżej jedna prosta, która przechodzi przez $A$ i jest równoległa do $l$. % lang-pl
    Per ogni retta $l$ e punto $A$, esiste al più una retta che passa per $A$ ed è parallela a $l$. % lang-it
\end{axiom}
\index{aksjomat!Playfaira}% % lang-pl
\index{assioma!di Playfair}% % lang-it
% TODO: https://en.wikipedia.org/wiki/Playfair%27s_axiom

John Playfair opublikuje ten aksjomat w 1795 roku, chociaż już wtedy stwierdzi, że inni używali go przed nim (np. William Ludlam). % lang-pl
John Playfair pubblicò questo assioma nel 1795, pur osservando che altri lo avevano già utilizzato prima di lui (ad esempio William Ludlam). % lang-it
\index[persons]{Playfair, John}%
\index[persons]{Ludlam, William}%

\begin{example}
    Rozważmy geometrię, gdzie płaszczyzna składa się z pięciu punktów $A$, $B$, $C$, $D$, $E$ i leży na niej dziesięć różnych prostych, każda z nich przechodząca przez dokładnie dwa różne punkty. % lang-pl
    Wtedy proste $AB$ i $AC$ mają punkt wspólny $A$, chociaż obydwie są równoległe do prostej $DE$. % lang-pl
    Aksjomat Playfaira nie jest spełniony. % lang-pl
    \index{aksjomat!Playfaira} % lang-pl
    Consideriamo una geometria in cui il piano è costituito da cinque punti $A$, $B$, $C$, $D$, $E$ e contiene dieci rette distinte, ciascuna passante per esattamente due punti distinti. % lang-it
    Allora le rette $AB$ e $AC$ hanno in comune il punto $A$, sebbene entrambe siano parallele alla retta $DE$. % lang-it
    L'assioma di Playfair non è soddisfatto. % lang-it
    \index{assioma!di Playfair} % lang-it
\end{example}

%

\subsection{Aksjomaty incydencji}
\index{aksjomat!incydencji}%

\begin{axiom}[incydencji, I1]
    Dla każdej pary punktów $A$ oraz $B$ istnieje dokładnie jedna prosta $l$, która przechodzi przez te punkty.
\end{axiom}

Zdania ,,prosta $l$ przechodzi przez punkt $A$'' oraz ,,punkt $A$ leży na prostej $l$'' mają dokładnie to samo znaczenie i będziemy je stosować zamiennie.

\begin{axiom}[incydencji, I2]
    Na każdej prostej leżą co najmniej dwa punkty.
\end{axiom}

\begin{axiom}[incydencji, I3]
    Istnieją co najmniej trzy punkty, które nie są współliniowe.
\end{axiom}

,,Być współliniowym'' to synonim leżenia na jednej prostej.
\index{współliniowy}%
Analogicznie mówimy, że wiele prostych, które przechodzą przez jeden punkt, jest współpękowych.

\begin{proposition}
    Dwie różne proste mogą mieć co najwyżej jeden punkt wspólny.
\end{proposition}

\begin{definition}
    Dwie różne proste, które nie są współpękowe (nie mają punktów wspólnych), nazywamy równoległymi.
    Każda prosta jest też równoległa do siebie.
\end{definition}
\index{prosta!równoległa}%

Symbolicznie: $k \parallel l$.
Trzeba uważać, bo dla niektórych dwie proste równoległe są współpękowe, bo przechodzą przez punkt w nieskończoności, jak zwykło się mawiać w geometrii rzutowej.

Na dwóch prostych, które nie są równoległe ($k \nparallel l$), leży dokładnie jeden punkt.
Mówimy, że proste się przecinają, a wspomniany punkt jest ich punktem przecięcia.
Oznaczamy go $k \cdot l$.

Hartshorne \cite[s. 38]{hartshorne_2000} przywoła:

\begin{axiom}[Playfaira, P]
    Dla każdej prostej $l$ oraz punktu $A$, istnieje co najwyżej jedna prosta, która przechodzi przez $A$ i jest równoległa do $l$.
\end{axiom}
\index{aksjomat!Playfaira}%
% TODO: https://en.wikipedia.org/wiki/Playfair%27s_axiom

John Playfair opublikuje ten aksjomat w 1795 roku, chociaż już wtedy stwierdzi, że inni używali go przed nim (np. William Ludlam).
\index[persons]{Playfair, John}%
\index[persons]{Ludlam, William}%

\begin{example}
    Rozważmy geometrię, gdzie płaszczyzna składa się z pięciu punktów $A$, $B$, $C$, $D$, $E$ i leży na niej dziesięć różnych prostych, każda z nich przechodząca przez dokładnie dwa różne punkty.
    Wtedy proste $AB$ i $AC$ mają punkt wspólny $A$, chociaż obydwie są równoległe do prostej $DE$.
    Aksjomat Playfaira nie jest spełniony. \index{aksjomat!Playfaira}
\end{example}

Aksjomat Playfaira jest równoważny z (piątym) postulatem Euklidesa i dlatego wiele osób spróbuje wyprowadzić go z czterech wcześniejszych postulatów. % lang-pl
Marek Kordos opisze w $\Delta_{12}^5$ (i $\Delta_{12}^6$) trzy ,,dowody'': Bolyaia, Legendre'a i Saccheriego. % lang-pl
L'assioma di Playfair è equivalente al (quinto) postulato di Euclide e per questo motivo molte persone tenteranno di dedurlo dai quattro postulati precedenti. % lang-it
Marek Kordos descriverà in $\Delta_{12}^5$ (e $\Delta_{12}^6$) tre ``dimostrazioni'': quelle di Bolyai, Legendre e Saccheri. % lang-it
\index[persons]{Bolyai, Farkas}%
\index[persons]{Legendre, Adrien-Marie}%
\index[persons]{Saccheri, Girolamo}%
% https://www.deltami.edu.pl/2012/05/dowody-v-postulatu-euklidesa/
Za każdym razem okaże się, że w ,,dowodzie'' użyte jest zdanie będące równoważnikiem piątego postulatu, na przykład: % lang-pl
Ogni volta si scoprirà che nella ``dimostrazione'' viene utilizzata un'affermazione equivalente al quinto postulato, ad esempio: % lang-it
\begin{itemize}
    \item suma kątów wewnętrznych pewnego trójkąta jest kątem półpełnym,\index{suma kątów}\index{somma degli angoli}% % lang-pl
    \item suma kątów wewnętrznych każdego trójkąta jest kątem półpełnym,\index{suma kątów}\index{somma degli angoli}% % lang-pl
    \item suma kątów wewnętrznych każdego trójkąta jest taka sama,\index{suma kątów}\index{somma degli angoli}% % lang-pl
    \item istnieją dwa trójkąty, które są do siebie podobne, ale nie przystające,\index{podobieństwo}\index{przystawanie}\index{similitudine}\index{congruenza}% % lang-pl
    \item na każdym trójkącie można opisać okrąg,\index{okrąg!opisany}\index{circonferenza!circoscritta}% % lang-pl
    \item istnieje para prostych, które są w stałej odległości od siebie, % lang-pl
    \item dwie proste, które są równoległe do danej prostej, są też równoległe do siebie,\index{prosta!równoległa}\index{retta!parallela}% % lang-pl
    \item twierdzenie Pitagorasa lub jego uogólnienie, twierdzenie kosinusów,\index{twierdzenie!Pitagorasa}\index{twierdzenie!kosinusów}\index{teorema!di Pitagora}\index{teorema!del coseno}% % lang-pl
    \item aksjomat Wallisa: nie ma górnego ograniczenia na pole trójkąta,\index{aksjomat!Wallisa}\index{assioma!di Wallis}% % lang-pl
    \item aksjomat Proklosa: jeśli prosta przecina jedną z dwóch równoległych prostych, to przecina także tę drugą,\index{aksjomat!Proklosa}\index{assioma!di Proclo}% % lang-pl
    \item kąty przy górnej podstawie czworokąta Saccheriego (czworokąta $ABCD$ o dwóch kątach prostych przy dolnej podstawie $AB$, w którym boki $AD$ i $BC$ mają równe długości) są proste,\index{czworokąt!Saccheriego}\index{quadrilatero!di Saccheri}% % lang-pl
    \item jeśli trzy kąty wewnętrzne czworokąta są proste, to czwarty kąt także jest prosty (czworokąt Lamberta jest prostokątem).\index{czworokąt!Lamberta}\index{quadrilatero!di Lambert}% % lang-pl
    \item la somma degli angoli interni di un certo triangolo è un angolo piatto,\index{suma kątów}\index{somma degli angoli}% % lang-it
    \item la somma degli angoli interni di ogni triangolo è un angolo piatto,\index{suma kątów}\index{somma degli angoli}% % lang-it
    \item la somma degli angoli interni di ogni triangolo è la stessa,\index{suma kątów}\index{somma degli angoli}% % lang-it
    \item esistono due triangoli simili ma non congruenti,\index{podobieństwo}\index{przystawanie}\index{similitudine}\index{congruenza}% % lang-it
    \item attorno a ogni triangolo si può circoscrivere una circonferenza,\index{okrąg!opisany}\index{circonferenza!circoscritta}% % lang-it
    \item esiste una coppia di rette che si trovano a distanza costante l'una dall'altra, % lang-it
    \item due rette parallele a una stessa retta sono parallele anche tra loro,\index{prosta!równoległa}\index{retta!parallela}% % lang-it
    \item il teorema di Pitagora oppure la sua generalizzazione, il teorema del coseno,\index{twierdzenie!Pitagorasa}\index{twierdzenie!kosinusów}\index{teorema!di Pitagora}\index{teorema!del coseno}% % lang-it
    \item assioma di Wallis: non esiste un limite superiore per l'area di un triangolo,\index{aksjomat!Wallisa}\index{assioma!di Wallis}% % lang-it
    \item assioma di Proclo: se una retta interseca una di due rette parallele, allora interseca anche l'altra,\index{aksjomat!Proklosa}\index{assioma!di Proclo}% % lang-it
    \item gli angoli alla base superiore del quadrilatero di Saccheri (il quadrilatero $ABCD$ con due angoli retti alla base inferiore $AB$ e con i lati $AD$ e $BC$ di uguale lunghezza) sono retti,\index{czworokąt!Saccheriego}\index{quadrilatero!di Saccheri}% % lang-it
    \item se tre angoli interni di un quadrilatero sono retti, allora anche il quarto è retto (il quadrilatero di Lambert è un rettangolo).\index{czworokąt!Lamberta}\index{quadrilatero!di Lambert}% % lang-it
% https://en.wikipedia.org/wiki/Saccheri_quadrilateral
\end{itemize}

% TODO: https://en.wikipedia.org/wiki/Parallel_postulate

\begin{proposition}
    Aksjomaty incydencji I1, I2, I3 oraz aksjomat P (Playfaira) są od siebie niezależne. % lang-pl
    \index{aksjomat!Playfaira} % lang-pl
    Gli assiomi di incidenza I1, I2, I3 e l'assioma P (di Playfair) sono indipendenti tra loro. % lang-it
    \index{assioma!di Playfair} % lang-it
\end{proposition}

Hartshorne \cite[s. 69-70]{hartshorne_2000} konstruuje modele geometrii, w których spełnione są dowolne trzy, ale nie czwarty z nich. % lang-pl
Hartshorne \cite[p. 69--70]{hartshorne_2000} costruisce modelli geometrici nei quali sono soddisfatti tre qualsiasi di essi, ma non il quarto. % lang-it

\subsubsection{Wstęp do geometrii rzutowej} % lang-pl
\subsubsection{Introduzione alla geometria proiettiva} % lang-it

\begin{proposition}
    Płaszczyzna rzutowa to taki zbiór punktów oraz prostych (podzbiorów zbioru punktów), że: % lang-pl
    \index{płaszczyzna!rzutowa} % lang-pl
    Un piano proiettivo è un insieme di punti e rette (sottoinsiemi dell'insieme dei punti) tale che: % lang-it
    \index{piano!proiettivo} % lang-it
    \begin{itemize}
        \item przez dwa różne punkty przechodzi dokładnie jedna prosta, % lang-pl
        \item każde dwie proste mają punkt wspólny, % lang-pl
        \item każda prosta ma co najmniej trzy punkty i % lang-pl
        \item nie wszystkie punkty są współliniowe. % lang-pl
        \item per due punti distinti passa una e una sola retta, % lang-it
        \item ogni due rette hanno un punto in comune, % lang-it
        \item ogni retta contiene almeno tre punti e % lang-it
        \item non tutti i punti sono allineati. % lang-it
    \end{itemize}
    (Każdy z wymienionych aksjomatów jest niezależny od pozostałych, ze wszystkich razem wynikają aksjomaty incydencji). % lang-pl
    Każda płaszczyzna rzutowa ma co najmniej 7 punktów, dokładnie jedna płaszczyzna rzutowa ma dokładnie 7 punktów. % lang-pl
    Jeśli istnieje prosta, która ma $n+1$ punktów, to płaszczyzna ma $n^2 + n + 1$ punktów. % lang-pl
    (Ciascuno degli assiomi elencati è indipendente dagli altri; tutti insieme implicano gli assiomi di incidenza). % lang-it
    Ogni piano proiettivo ha almeno 7 punti ed esiste un unico piano proiettivo con esattamente 7 punti. % lang-it
    Se esiste una retta che contiene $n+1$ punti, allora il piano contiene $n^2+n+1$ punti. % lang-it
\end{proposition}

Marek Kordos opisze dziewięć twarzy płaszczyzny rzutowej w $\Delta_{13}^5$, wcześniej czymś podobnym zajmie się Maria Donten-Bury w $\Delta_{11}^6$. % lang-pl
Marek Kordos descriverà i nove volti del piano proiettivo in $\Delta_{13}^5$; in precedenza Maria Donten-Bury si occuperà di un argomento simile in $\Delta_{11}^6$. % lang-it
% https://www.deltami.edu.pl/2013/05/dziewiec-twarzy-plaszczyzny-rzutowej/ % TODO
Temat będzie wracać: w $\Delta_{19}^4$ podejmie go Michał Krych. % lang-pl
L'argomento tornerà più volte: in $\Delta_{19}^4$ sarà ripreso da Michał Krych. % lang-it
% https://www.deltami.edu.pl/2019/04/co-to-jest-geometria-rzutowa/

\begin{example}
[płaszczyzna Fana] % lang-pl
[piano di Fano] % lang-it
    Zbiór złożony z siedmiu elementów (,,punktów''), w którym wyróżniono rodzinę siedmiu podzbiorów (,,prostych'') spełniający następujące warunki: % lang-pl
    \index{płaszczyzna!Fana} % lang-pl
    Chiamiamo piano di Fano un insieme costituito da sette elementi (``punti''), nel quale è stata distinta una famiglia di sette sottoinsiemi (``rette'') che soddisfano le seguenti condizioni: % lang-it
    \index{piano!di Fano} % lang-it
    \begin{itemize}
        \item każde dwie różne proste mają dokładnie jeden punkt wspólny, % lang-pl
        \item każde dwa różne punkty należą do dokładnie jednej prostej % lang-pl
        \item ogni due rette distinte hanno esattamente un punto in comune, % lang-it
        \item ogni due punti distinti appartengono a una e una sola retta % lang-it
    \end{itemize}
    nazywamy płaszczyzną Fana. % lang-pl
    Jest przykładem płaszczyzny rzutowej, która nie spełnia aksjomatu Fana (że trzy punkty przekątne czworoboku zupełnego nie są współliniowe). % lang-pl
    \index{aksjomat!Fana}5 % lang-pl
    È un esempio di piano proiettivo che non soddisfa l'assioma di Fano (secondo il quale i tre punti diagonali di un quadrilatero completo non sono allineati). % lang-it
    \index{assioma!di Fano}5 % lang-it
    \begin{figure}[H]
        \centering
        \begin{tikzpicture}[
        mydot/.style={
        draw,
        circle,
        fill=black,
        inner sep=1.5pt}
        ]
        \draw
        (0,0) coordinate (A) --
        (4cm,0) coordinate (B) --
        ($ (A)!.5!(B) ! {sin(60)*2} ! 90:(B) $) coordinate (C) -- cycle;
        \coordinate (O) at
        (barycentric cs:A=1,B=1,C=1);
        \draw (O) circle [radius=4cm*1.717/6];
        \draw (C) -- ($ (A)!.5!(B) $) coordinate (LC);
        \draw (A) -- ($ (B)!.5!(C) $) coordinate (LA);
        \draw (B) -- ($ (C)!.5!(A) $) coordinate (LB);
        \foreach \Nodo in {A,B,C,O,LC,LA,LB}
        \node[mydot] at (\Nodo) {};
    \end{tikzpicture}%
    \caption{Płaszczyzna Fana z 1892 roku} % lang-pl
    \caption{Piano di Fano del 1892} % lang-it
\end{figure}
\end{example}

Jest po jednej płaszczyźnie rzutowej dla $n = 3, 4, 5, 6, 8$ punktów na każdej prostej; żadnej dla $n = 7$, co najmniej jedna dla $n = 9$ i co najmniej cztery dla $n = 10$. % lang-pl
O płaszczyźnie Fana napisze Eves \cite[s. 366]{eves1_1972}. % lang-pl
Esiste un unico piano proiettivo con $n = 3, 4, 5, 6, 8$ punti su ogni retta; nessuno per $n = 7$, almeno uno per $n = 9$ e almeno quattro per $n = 10$. % lang-it
Del piano di Fano scriverà Eves \cite[p. 366]{eves1_1972}. % lang-it

W momencie pisania tego tekstu nie interesuje nas jeszcze geometria rzutowa, więc wymienimy tylko kilka terminów: % lang-pl
dwustosunek, % lang-pl
\index{dwustosunek}% % lang-pl
rzutowe przekształcenia łańcuchów, stożkowych oraz pęków stycznych do stożkowych. % lang-pl
\index{stożkowa}% % lang-pl
Oprócz tego są dwa twierdzenia: Steinera i Braikendridge'a-Maclaurina. % lang-pl
\index{twierdzenie!Steinera}% % lang-pl
\index{twierdzenie!Braikendridge'a-Maclaurina}% % lang-pl
Al momento della stesura di questo testo non ci interessa ancora la geometria proiettiva, quindi citeremo soltanto alcuni termini: % lang-it
birapporto, % lang-it
\index{birapporto}% % lang-it
trasformazioni proiettive di catene, coniche e fasci tangenti a coniche. % lang-it
\index{conica}% % lang-it
Inoltre vi sono due teoremi: quello di Steiner e quello di Braikenridge--Maclaurin. % lang-it
\index{teorema!di Steiner}% % lang-it
\index{teorema!di Braikenridge-Maclaurin}% % lang-it

Patrz też do Audina \cite[s. 143-182]{audin_2003}. % lang-pl
Si veda anche Audin \cite[p. 143--182]{audin_2003}. % lang-it
% Delta: 2024/lipiec.

\subsubsection{Wstęp do geometrii afinicznej} % lang-pl
\subsubsection{Introduzione alla geometria affine} % lang-it

\begin{proposition}
    Płaszczyzna afiniczna to taki zbiór punktów i prostych, które spełniają: % lang-pl
    \index{płaszczyzna!afiniczna} % lang-pl
    Un piano affine è un insieme di punti e rette che soddisfa: % lang-it
    \index{piano!affine} % lang-it
    \begin{itemize}
        \item aksjomaty incydencji oraz % lang-pl
        \item mocniejszą wersję aksjomatu Playfaira: dla każdej prostej $l$ i punktu $A$, dokładnie jedna prosta przechodzi przez punkt $A$ i jest równoległa do $l$. % lang-pl
        \item gli assiomi di incidenza e % lang-it
        \item una versione più forte dell'assioma di Playfair: per ogni retta $l$ e punto $A$, passa per $A$ una e una sola retta parallela a $l$. % lang-it
    \end{itemize}
    Każda prosta na płaszczyźnie afinicznej ma tyle samo punktów. % lang-pl
    Jeśli pewna prosta ma $n$ punktów, to płaszczyzna ma dokładnie $n^2$ punktów. % lang-pl
    Istnieją płaszczyzny afiniczne o $4$, $9$, $16$ i $25$ punktach, ale nie istnieje taka, która miałaby $36$ punktów. % lang-pl
    Ogni retta di un piano affine contiene lo stesso numero di punti. % lang-it
    Se una retta contiene $n$ punti, allora il piano contiene esattamente $n^2$ punti. % lang-it
    Esistono piani affini con $4$, $9$, $16$ e $25$ punti, ma non ne esiste alcuno con $36$ punti. % lang-it
\end{proposition} % Hartshorne 71, 72

Podręcznik Audina \cite[s. 7]{audin_2003} zaczyna się od definicji przestrzeni afinicznej. % lang-pl
Natomiast Coxeter \cite[s. 209]{coxeter_1967} wspomni, że według Blaschkego słowo ,,afiniczny'' utworzył Euler. % lang-pl
Powinowactwo (bo tak tłumaczy się na polski \emph{,,affinity''}) znajdzie się w niektórych numerach czasopisma Delta: $\Delta_{24}^6$, ale nie za bardzo u nas. % lang-pl
Il manuale di Audin \cite[p. 7]{audin_2003} inizia con la definizione di spazio affine. % lang-it
Coxeter \cite[p. 209]{coxeter_1967} ricorda invece che, secondo Blaschke, il termine ``affine'' fu coniato da Euler. % lang-it
L'affinità (che è il significato del termine inglese \emph{``affinity''}) compare in alcuni numeri della rivista Delta, ad esempio in $\Delta_{24}^6$, ma non avrà grande spazio in questo testo. % lang-it

%